% Generated by Listas/make_statements.py from pratica-lecture-01.md, item 16. Do not edit.
\textbf{Source.} Tadelis (2013), \emph{Game Theory: An Introduction}, Exercise 4.7 ``Campaigning'', pp.~77--78 {[}PDF 96--97{]}. Reproduced verbatim. Answer in the source: none in the book (the preface, p.~xiii, says about half are solved on Princeton UP's website; not checked).



\begin{quote}

\textbf{Campaigning:} Two candidates, 1 and 2, are running for office. Each has one of three choices in running his campaign: focus on the positive aspects of one's own platform (call this a positive campaign {[}\(P\){]}), focus on the positive aspects of one's own platform while attacking one's opponent's campaign (call this a balanced campaign {[}\(B\){]}), or finally focus only on attacking one's opponent (call this a negative campaign {[}\(N\){]}). All a candidate cares about is the probability of winning, so assume that if a candidate expects to win with probability \(\pi \in [0, 1]\), then his payoff is \(\pi\). The probability that a candidate wins depends on his choice of campaign and his opponent's choice. The probabilities of winning are given as follows:



If both choose the same campaign then each wins with probability 0.5.



If candidate \(i\) uses a positive campaign while \(j \neq i\) uses a balanced one then \(i\) loses for sure.



If candidate \(i\) uses a positive campaign while \(j \neq i\) uses a negative one then \(i\) wins with probability 0.3.



If candidate \(i\) uses a negative campaign while \(j \neq i\) uses a balanced one then \(i\) wins with probability 0.6.



\begin{enumerate}

\def\labelenumi{\alph{enumi}.}

\item

  Model this story as a normal-form game. (It suffices to be specific about the payoff function of one player and to explain how the other player's payoff function is different and why.)

\item

  Write down the game in matrix form.

\item

  What happens at each stage of elimination of strictly dominated strategies? Will this procedure lead to a clear prediction?

\end{enumerate}

\end{quote}

